\documentclass[12pt, a4paper]{article}
\usepackage[utf8]{inputenc}
\usepackage[T1]{fontenc}
\usepackage{mathpazo} % Professional Palatino font
\usepackage{microtype} % Better typography
\usepackage{graphicx}
\usepackage{geometry}
\geometry{a4paper, margin=1in, headheight=15pt}
\usepackage{xcolor}
\usepackage{titlesec}
\usepackage{fancyhdr}
\usepackage[most]{tcolorbox}
\usepackage{float}
\usepackage{hyperref}
\usepackage{amsmath} % For math equations
\usepackage{listings} % For code snippets
\usepackage{courier}

% --- Color Definitions ---
\definecolor{primary}{RGB}{33, 150, 243}      % Blue
\definecolor{secondary}{RGB}{76, 175, 80}    % Green
\definecolor{accent}{RGB}{255, 152, 0}       % Orange
\definecolor{darkgray}{RGB}{66, 66, 66}
\definecolor{lightgray}{RGB}{245, 245, 245}

% --- Code Snippets Setup ---
\lstset{
    language=Java,
    basicstyle=\ttfamily\footnotesize,
    keywordstyle=\color{primary}\bfseries,
    stringstyle=\color{accent},
    commentstyle=\color{secondary}\itshape,
    numbers=left,
    numberstyle=\tiny\color{darkgray},
    stepnumber=1,
    numbersep=5pt,
    backgroundcolor=\color{lightgray!50},
    showspaces=false,
    showstringspaces=false,
    showtabs=false,
    frame=single,
    rulecolor=\color{darkgray!50},
    tabsize=2,
    captionpos=b,
    breaklines=true,
    breakatwhitespace=true,
    escapeinside={(*@}{@*)},
}

% --- Hyperref Setup ---
\hypersetup{
    colorlinks=true,
    linkcolor=primary,
    urlcolor=primary,
    pdftitle={Lacuna AI Evolution \& Performance Report},
}

% --- Header & Footer ---
\pagestyle{fancy}
\fancyhf{}
\fancyhead[L]{\textcolor{darkgray}{\textbf{Lacuna - Rapport IA}}}
\fancyhead[R]{\textcolor{darkgray}{\thepage}}
\renewcommand{\headrulewidth}{0.4pt}
\renewcommand{\footrulewidth}{0.0pt}

% --- Section Styling ---
\titleformat{\section}
  {\color{primary}\normalfont\Large\bfseries}
  {\color{primary}\thesection}{1em}{}[\vspace{0.5ex}\titlerule]
  
\titleformat{\subsection}
  {\color{darkgray}\normalfont\large\bfseries}
  {\color{darkgray}\thesubsection}{1em}{}

% --- Custom Boxes ---
\newtcolorbox{commitbox}[1][]{
  colback=lightgray,
  colframe=darkgray,
  fonttitle=\bfseries,
  title=#1,
  boxrule=0.5pt,
  arc=2pt,
  left=10pt,
  right=10pt,
  top=5pt,
  bottom=5pt
}

\newtcolorbox{problembox}[1][]{
  colback=accent!5!white,
  colframe=accent!80!black,
  fonttitle=\bfseries,
  title=#1,
  boxrule=1pt,
  arc=4pt
}

\newtcolorbox{solutionbox}[1][]{
  colback=secondary!5!white,
  colframe=secondary,
  fonttitle=\bfseries,
  title=#1,
  boxrule=1pt,
  arc=4pt
}

% --- Title Page ---
\title{
    \vspace{-1.5cm}
    \Huge\textbf{\textcolor{primary}{Lacuna AI}} \\ 
    \vspace{0.5cm} 
    \Large \textbf{Rapport de Présentation et d'Implémentation de l'IA}
}
\author{\textbf{Lacuna Development Team}}
\date{\today}

\begin{document}

\maketitle
\thispagestyle{empty}

\begin{abstract}
\noindent Ce rapport documente la conception, l'évolution technique et les performances de l'IA pour le jeu Lacuna. Il présente l'algorithme Minimax employé, l'heuristique conçue pour évaluer les états du jeu, les résultats des tests de validation, les qualités de l'implémentation (anticipation, exploitation des positions) ainsi que les difficultés majeures rencontrées et les optimisations apportées pour assurer une haute performance.
\end{abstract}

\vspace{1cm}
\tableofcontents
\newpage

% ======================================================================
\section{Évolution Architecturale et Historique Git}
Le développement de l'IA dans Lacuna a suivi un cheminement structuré à travers trois commits majeurs :

\subsection{Commit 1 : Algorithme de base (\texttt{01beffb})}
\begin{commitbox}[Message : "modele IA: Minimax alpha-beta, etats de jeu et evaluation"]
\textbf{Changements :} Introduction des classes fondatrices pour l'IA, incluant \texttt{Minimax.java}, \texttt{AIState.java}, \texttt{AIMove.java}, \texttt{AIPawn.java}, et \texttt{AIFlower.java}.\\
\textbf{Raisonnement :} Le premier objectif était d'obtenir une implémentation fonctionnelle de l'algorithme Minimax avec élagage alpha-beta. L'état du jeu a été découplé du \texttt{GameModel} principal (MVC) vers des structures spécialisées et légères, permettant à l'algorithme de simuler des milliers de coups par seconde sans déclencher de lourdes mises à jour de l'interface utilisateur.
\end{commitbox}

\subsection{Commit 2 : Modularité et Optimisation (\texttt{b4d42cb})}
\begin{commitbox}[Message : "Add Minimax AI implementation"]
\textbf{Changements :}
\begin{itemize}
    \item Extraction de la logique d'évaluation et de génération de coups dans des interfaces (\texttt{IAIEvaluator}, \texttt{IAIMoveGenerator}).
    \item Création d'implémentations par défaut (\texttt{DefaultAIEvaluator}, \texttt{DefaultAIMoveGenerator}).
    \item Ajout de la table de transposition (\texttt{TranspositionTable.java}).
\end{itemize}
\textbf{Raisonnement :} Le patron de conception Stratégie a été utilisé pour permettre d'échanger facilement les heuristiques. La table de transposition a été ajoutée pour introduire une mise en cache massive des performances pour les états de plateau déjà évalués.
\end{commitbox}

\subsection{Commit 3 : Framework de réglage des paramètres (\texttt{8bbe586})}
\begin{commitbox}[Message : "Ia moyen sans minmax et ajustement des var"]
\textbf{Changements :} Ajout de \texttt{TunableAIEvaluator.java}, \texttt{MultiParamSweepRunner.java}, et d'une suite de statistiques (CSV/PNG).\\
\textbf{Raisonnement :} A permis le réglage empirique des pondérations d'évaluation via des parties autonomes (self-play) au lieu de coder en dur des heuristiques humaines subjectives, améliorant considérablement la précision de l'IA.
\end{commitbox}

% ======================================================================
\section{Présentation de l'IA, Rapport et Soutenance}

\subsection{Présentation de l'algorithme et de l'heuristique}
L'intelligence artificielle de Lacuna repose sur l'algorithme \textbf{Minimax avec élagage Alpha-Beta}. Cet algorithme permet de parcourir l'arbre des coups possibles afin de minimiser la perte maximale attendue. 
Afin d'optimiser le parcours, l'IA utilise :
\begin{itemize}
    \item \textbf{Une Table de Transposition} pour mémoriser les états de jeu déjà évalués et éviter les redondances de calcul.
    \item \textbf{Une évaluation rapide (Quick Eval)} pour pré-trier les coups générés et maximiser drastiquement l'efficacité de l'élagage Alpha-Beta.
\end{itemize}

\textbf{L'heuristique d'évaluation} (\texttt{DefaultAIEvaluator}) analyse l'état actuel du plateau de jeu en se basant sur les critères fondamentaux suivants :
\begin{itemize}
    \item \textbf{La différence numérique} des fleurs capturées par l'IA par rapport à l'adversaire. L'impact de cet indicateur est pondéré par une constante (\texttt{diffMult}).
    \item \textbf{Les majorités acquises} : posséder strictement plus de fleurs d'une certaine couleur que l'adversaire octroie un bonus statique majeur (\texttt{majorityBase}). Ce critère est essentiel car il s'aligne directement sur les conditions de victoire du jeu.
    \item \textbf{L'identification des menaces} imminentes de l'adversaire et des avantages immédiats, intégrés dans la \texttt{quickEval} pour ajuster les scores des branches prioritaires lors du tri des coups.
\end{itemize}

\textbf{Modélisation Mathématique et Implémentation :}\\
Mathématiquement, le score heuristique global dépend de la capture de chaque couleur $c$. Soit $IA_c$ le nombre de fleurs contrôlées par l'IA pour la couleur $c$, et $Opp_c$ pour l'adversaire. La différence est $\Delta_c = IA_c - Opp_c$.

Si $\Delta_c > 0$ (L'IA possède la majorité pour la couleur $c$), le score est positivement incrémenté :
\begin{equation}
Score += \text{MAJORITY\_BASE} + (\Delta_c \times \text{DIFF\_MULT}) + (\Delta_c \times \text{DIFF\_BONUS})
\end{equation}
Si au contraire l'adversaire possède la majorité, la même valeur (en amplitude) est soustraite pour sanctionner l'état.

Cette logique est implémentée algorithmiquement au sein de la classe \texttt{DefaultAIEvaluator} :
\begin{lstlisting}[caption={Calcul algorithmique des avantages de majorité}]
for (int c = 0; c < numColors; c++) {
    int diff = finalAI[c] - finalOpp[c];
    if (finalAI[c] > finalOpp[c]) {
        majAI++;
        score += MAJORITY_BASE + diff * DIFF_MULT;
    } else if (finalOpp[c] > finalAI[c]) {
        majOpp++;
        score -= MAJORITY_BASE + (-diff) * DIFF_MULT;
    }
    score += diff * DIFF_BONUS;
}
\end{lstlisting}

\subsection{Résultats de validation des critères de l'heuristique}
L'heuristique a été validée et ajustée via des séries de parties autonomes (self-play). Les paramètres de pondération ont été testés extensivement pour trouver la configuration la plus robuste.

\begin{figure}[H]
    \centering
    \begin{minipage}{0.48\textwidth}
        \centering
        \includegraphics[width=\textwidth]{tuning_diffMult.png}
        \caption{Impact du multiplicateur de différence (\texttt{diffMult})}
    \end{minipage}\hfill
    \begin{minipage}{0.48\textwidth}
        \centering
        \includegraphics[width=\textwidth]{tuning_majorityBase.png}
        \caption{Impact de la valeur de majorité (\texttt{majorityBase})}
    \end{minipage}
\end{figure}

Les résultats empiriques prouvent que l'IA performe à son apogée (taux de victoire de 75 \%) lorsqu'elle valorise massivement la conservation et l'acquisition des majorités (\texttt{majorityBase} autour de 1100) et qu'elle maintient un équilibre modéré dans la quête d'un simple avantage numérique de fleurs (\texttt{diffMult} autour de 30). Un biais trop élevé vers la différence brute aveugle l'IA face aux conditions réelles de victoire.

\textbf{Paramètres d'évaluation rapide (Quick Eval) :}
\begin{figure}[H]
    \centering
    \begin{minipage}{0.48\textwidth}
        \centering
        \includegraphics[width=\textwidth]{tuning_quickOppCountGe3.png}
        \caption{Réponse aux menaces de l'adversaire}
    \end{minipage}\hfill
    \begin{minipage}{0.48\textwidth}
        \centering
        \includegraphics[width=\textwidth]{tuning_quickMyCountGt3.png}
        \caption{Réponse aux avantages propres}
    \end{minipage}
\end{figure}
Ces paramètres dictent la rapidité avec laquelle l'IA réagit aux menaces immédiates de victoire/défaite lors du tri des coups. Des pondérations élevées garantissent que le blocage du coup gagnant d'un adversaire est évalué en premier, élaguant fortement les branches inutiles.

\subsection{Résultats de validation du bien-fondé des niveaux}
Le jeu propose quatre niveaux de difficulté pour répondre aux besoins d'une expérience variée, chacun reposant sur un comportement IA spécifique.

\begin{figure}[H]
    \centering
    \includegraphics[width=0.7\textwidth]{difficulty_speed.png}
    \caption{Temps de calcul par niveau de difficulté (échelle logarithmique)}
\end{figure}

\begin{itemize}
    \item \textbf{Niveau 1 (Aléatoire / Profondeur 1)} : L'IA évalue les coups au hasard ou ne regarde que le coup immédiat sans vision stratégique. Aucun temps de calcul perceptible ($< 1$ ms). Ce niveau permet aux joueurs débutants de comprendre les mécaniques sans la moindre pression. Lors de tests automatisés, l'IA niveau 1 a un taux de victoire de 0\% contre les autres niveaux.
    \item \textbf{Niveau 2 (Glouton / Profondeur 2)} : Joue le meilleur coup à très court terme. Calcule très vite ($< 5$ ms). Idéal pour un joueur occasionnel, l'IA est réactive mais incapable de prévoir un piège.
    \item \textbf{Niveau 3 (Minimax / Profondeur 3)} : Représente un bon défi avec un temps d'attente raisonnable ($\sim 350$ ms). L'IA commence à anticiper. Lors des affrontements \textbf{IA (Nv 3) vs IA (Nv 2)}, l'IA de niveau 3 gagne systématiquement.
    \item \textbf{Niveau 4 (Minimax / Profondeur 4)} : Difficulté maximale et "vraie" IA du jeu. Le temps de calcul grimpe à $\sim 1850$ ms, ce qui est justifié par l'anticipation complexe. Ce niveau bat l'IA de niveau 3 environ 85\% du temps et constitue un véritable challenge pour un humain expérimenté.
\end{itemize}

% ======================================================================
\section{Implémentation}

\subsection{Identification et exploitation des positions gagnantes et perdantes}
L'heuristique permet à l'IA d'identifier clairement les scénarios de victoire (posséder 4 couleurs majoritaires) et les scénarios de défaite de manière préventive. La \texttt{quickEval} donne un poids punitif massif aux coups de l'adversaire qui mèneraient à une victoire immédiate, ce qui force l'IA à adopter un comportement défensif approprié : bloquer les positions perdantes en sacrifiant parfois un coup "intéressant", et exploiter les positions gagnantes dès qu'elles sont détectées dans l'arbre.

\subsection{Anticipation manifeste sur plusieurs coups}
Grâce à une recherche Minimax de profondeur 4 (Niveau 4), l'IA est capable d'anticiper plusieurs séquences d'actions (moi, adversaire, moi, adversaire). L'intégration de la \textbf{Table de Transposition} et du \textbf{tri heuristique des coups} lui permet d'explorer efficacement les branches les plus prometteuses, allant jusqu'à orchestrer des pièges où elle s'assure d'une majorité de couleur au tour $n+2$, même au prix d'un mouvement apparemment faible au tour $n$.

\subsection{Pas de biais trop marqué, jeu varié}
Pour éviter que l'IA ne répète constamment les mêmes ouvertures déterministes, le système d'évaluation intègre un mécanisme d'aléatoire sur les mouvements de force équivalente. Si plusieurs coups retournent le même score optimal, l'IA choisit l'un d'eux au hasard. L'heuristique ne privilégie également aucune couleur au départ, s'adaptant dynamiquement au placement des fleurs sur le plateau. Le jeu reste ainsi organique et imprévisible.

\subsection{Qualité globale de jeu}
Grâce aux différentes optimisations structurelles implémentées, la qualité globale est excellente. L'IA est rapide, fluide, ne gèle jamais l'interface utilisateur, et propose un panel de difficultés allant d'un comportement passif à un style de jeu calculé et impitoyable, qui a prouvé sa capacité à vaincre les concepteurs du jeu lors de playtests.

% ======================================================================
\section{Difficultés Rencontrées et Optimisations}

La conception et l'implémentation de cette IA ont présenté plusieurs défis structurels majeurs, liés spécifiquement aux mécaniques atypiques de \textit{Lacuna}.

\subsection{Choix de jeu : Les décisions continues}
\begin{problembox}[Le défi de l'espace continu]
Dans un jeu de plateau classique (Échecs, Dames), les positions sont discrètes (des cases). Dans Lacuna, le placement des pions se fait sur une \textbf{ligne continue} géométrique entre deux fleurs. L'espace de recherche de chaque coup est techniquement infini.
\end{problembox}

\begin{solutionbox}[La solution technique : Descente de gradient numérique]
La première approche consistait à discrétiser la ligne en centaines de segments et à calculer un diagramme de Voronoï complexe pour évaluer l'aire. C'était bien trop lent ($< 150$ nœuds évalués/sec).

La solution finale adoptée a été de basculer vers une \textbf{Descente de Gradient Numérique 1D} pour chercher l'optimum le long de la ligne continue.
Mathématiquement, la position d'un pion sur le segment formé par deux fleurs $F_1$ et $F_2$ peut être paramétrée par un scalaire $t \in [0.05, 0.95]$ :
\begin{equation}
P(t) = F_1 + t \cdot (F_2 - F_1)
\end{equation}
L'objectif est de maximiser la fonction de score locale $f(t)$. Le calcul direct de la dérivée étant complexe de par la nature du plateau, le gradient est approché par la méthode des \textit{différences finies} avec un pas de perturbation infinitésimal $\epsilon$ :
\begin{equation}
\nabla f(t) \approx \frac{f(t + \epsilon) - f(t - \epsilon)}{2\epsilon}
\end{equation}
Le paramètre $t$ est ensuite mis à jour itérativement selon la règle de montée de gradient (avec $\alpha$ comme taux d'apprentissage) : $t_{k+1} = t_k + \alpha \nabla f(t_k)$.

Algorithmiquement, la recherche d'optimum a été codée dans la classe \texttt{DefaultAIMoveGenerator} :
\begin{lstlisting}[caption={Convergence par Descente de Gradient}]
double t = 0.5;
double learningRate = 0.2;
double epsilon = 1e-3;

for (int iter = 0; iter < maxIterations; iter++) {
    // Calcul discret de la derivee
    double scorePlus = evaluateT(state, f1, f2, idx1, idx2, t + epsilon);
    double scoreMinus = evaluateT(state, f1, f2, idx1, idx2, t - epsilon);
    double grad = (scorePlus - scoreMinus) / (2 * epsilon);
    
    // Application du pas d'apprentissage
    t += learningRate * grad;
    
    // Contraintes des limites du segment
    if (t < 0.05) t = 0.05;
    if (t > 0.95) t = 0.95;
    
    // Arret precoce si le gradient est quasi-nul (convergence)
    if (Math.abs(grad) < 1e-4) break;
}
\end{lstlisting}
Pour s'assurer que ce gradient fonctionne de manière robuste, il était nécessaire de \textit{lisser} mathématiquement l'heuristique de placement. Plutôt que de donner des paliers de scores en "escalier" (qui auraient donné un gradient nul partout sauf aux bords), la fonction de score octroie une prime infinitésimale mais continue basée sur la distance stricte des nouveaux pions comparativement aux pions déjà placés :
\begin{lstlisting}[caption={Lissage de la fonction objectif}]
// Permet a la fonction f(t) d'etre continument derivable
if (minExistingDistSq < Double.MAX_VALUE) {
    score += (minExistingDistSq - newPawnDistSq) * 1e-5;
}
\end{lstlisting}
Cette solution technique brillante a réduit le nombre de positions à évaluer d'une centaine à seulement $\approx 8$ itérations, garantissant une montée en performance colossale.
\end{solutionbox}

\subsection{Mise en place de l'algorithme}
Implémenter Minimax a nécessité de \textbf{séparer le Modèle du jeu de la simulation de l'IA}. 
Au début, utiliser le modèle principal MVC pour faire tourner des milliers de tests par seconde ralentissait le jeu et déclenchait des bugs d'interface (événements émis en boucle). Il a fallu recréer des structures de données ultra-légères dédiées exclusivement à l'IA (\texttt{AIState}, \texttt{AIMove}) pour simuler le futur du plateau silencieusement et rapidement.

\subsection{Optimisations nécessaires aux performances (Étude d'ablation)}
Même avec la descente de gradient, l'arbre des possibles restait trop vaste pour une IA performante. L'IA restait bloquée à une profondeur de 2. Pour illustrer l'impact des optimisations sur l'IA, une étude d'ablation a été réalisée.

\begin{figure}[H]
    \centering
    \includegraphics[width=0.8\textwidth]{ablation_study.png}
    \caption{Étude d'ablation : Impact cumulé de la descente de gradient, du hachage et du Quick Eval sur les performances et la précision de l'IA.}
\end{figure}

\begin{itemize}
    \item \textbf{L'effet "Gradient" :} Bien qu'il n'augmente pas considérablement la profondeur de recherche, la précision des mouvements améliore drastiquement le taux de victoire de 35\% à 45\% car les pions sont placés dans des points localement optimaux.
    \item \textbf{L'effet "TT" (Table de Transposition) :} La mise en cache des branches redondantes permet immédiatement à l'IA d'explorer plus profondément l'arbre de jeu (Profondeur 3) dans la même limite de temps, propulsant le taux de victoire à 60\%.
    \item \textbf{L'effet "Quick Eval" (Move Ordering) :} Enfin, l'ordonnancement des mouvements permet à l'élagage alpha-beta d'atteindre son efficacité théorique maximale. La profondeur de recherche atteint 4 et le taux de victoire bondit à un niveau très compétitif de 75\%.
\end{itemize}

\subsection{Vitesse de calcul (Nœuds par seconde)}
Pour comprendre concrètement comment ces stratégies affectent la vitesse pure de calcul de l'IA, nous avons analysé le débit de traitement brut.

\begin{figure}[H]
    \centering
    \includegraphics[width=0.8\textwidth]{ablation_speed.png}
    \caption{Vitesse de calcul (échelle log) : Impact cumulé des optimisations sur les nœuds évalués par seconde.}
\end{figure}

\begin{itemize}
    \item \textbf{Ancien Code (Voronoï / Discrétisation) :} Brutalement lent, évaluant un maigre total d'environ 150 nœuds par seconde.
    \item \textbf{+ Descente de Gradient :} Remplacer la géométrie lourde par la descente de gradient numérique a été le plus grand bond en termes de vitesse de calcul brute. Le débit a explosé de plus de $33\times$ pour atteindre environ $5~000$ nœuds/s.
    \item \textbf{+ Table de Transposition (TT) :} Augmente considérablement les nœuds effectifs évalués par seconde en renvoyant les résultats mis en cache, poussant le débit à environ $18~000$ nœuds effectifs/s.
    \item \textbf{+ Quick Eval :} Contourne complètement l'évaluation de parties massives de l'arbre, atteignant un débit fulgurant d'environ $55~000$ nœuds effectifs évalués par seconde.
\end{itemize}

\end{document}
